% !TEX TS-program = lualatex
\documentclass{article}
\usepackage[margin=2cm]{geometry}
\usepackage{tkz-fn}
\usetikzlibrary{arrows}
\parindent=0pt
\begin{document}

\section*{1. Polynôme, point, tangente (format us)}
\tkzFnSetup{format=us}
\tkzFnDef{f}{x^2 - 2*x - 1}
\begin{tikzpicture}
  \tkzFnInit[xmin=-2,xmax=4,ymin=-3,ymax=6]
  \tkzFnGrid
  \tkzFnAxes
  \begin{scope}
    \tkzFnClip
    \tkzFnDraw[thick,red](f)
  \end{scope}
  \tkzFnDefPoint(f,3){A}
  \tkzFnTangent[blue,kl=1.5,kr=1.5](f,3)
  \tkzFnDrawPoints[red](A)
  \node[above left] at (A) {$A$};
\end{tikzpicture}

$f(3)=\tkzFnPrint{f(3)}$, $f'(3)\approx\tkzFnPrint[4]{diff(f,3)}$,
$A\tkzFnPrintCoord(A)$, $f(0.25)=\tkzFnPrint[4]{f(0.25)}$,
$\tkzFnPrint{1234567.891}$, $I=\tkzFnPrintInterval[oc]{-\infty}{2.5}$

\section*{2. Mêmes résultats en format fr}
\tkzFnSetup{format=fr}
$f(3)=\tkzFnPrint{f(3)}$, $f'(3)\approx\tkzFnPrint[4]{diff(f,3)}$,
$A\tkzFnPrintCoord(A)$, $f(0{,}25)=\tkzFnPrint[4]{f(0.25)}$,
$\tkzFnPrint{1234567.891}$, $I=\tkzFnPrintInterval[oc]{-\infty}{2.5}$,
$J=\tkzFnPrintInterval[oo]{0}{pi}$

\section*{3. Asymptotes : $1/x$, $\tan x$, $\ln x$ (format fr)}
\begin{tikzpicture}[scale=.8]
  \tkzFnInit[xmin=-4,xmax=4,ymin=-4,ymax=4]
  \tkzFnGrid
  \tkzFnAxes
  \begin{scope}
    \tkzFnClip
    \tkzFnDraw[thick,blue,samples=400]{1/x}
    \tkzFnDraw[thick,orange,samples=600]{tan(x)}
    \tkzFnDraw[thick,green!50!black,domain=0:4]{ln(x)}
  \end{scope}
\end{tikzpicture}

\section*{4. Unités : $x\mapsto x^3$ sur $[0;100]$, xstep=10, ystep=100000}
\begin{tikzpicture}[scale=.5]
  \tkzFnInit[xmin=0,xmax=100,xstep=10,ymin=0,ymax=1000000,ystep=100000]
  \tkzFnGrid
  \tkzFnAxes[font=\scriptsize]
  \tkzFnDef{g}{x^3}
  \tkzFnDraw[red,thick](g)
  \tkzFnDefPoint(g,50){B}
  \tkzFnTangent[blue,kl=20,kr=20](g,50)
  \tkzFnDrawPoints(B)
\end{tikzpicture}

$B\tkzFnPrintCoord(B)$, pente en $B$ : $\tkzFnPrint{diff(g,50)}$
(dans le repère, et non en cm).

\section*{5. Paramètre Lua et fonction composée}
\directlua{a = 0.5}
\tkzFnDef{h}{a*f(x) + 1}
\begin{tikzpicture}[scale=.6]
  \tkzFnInit[xmin=-3,xmax=5,ymin=-2,ymax=6,xorigin=0]
  \tkzFnAxes[labels=false]
  \tkzFnDraw[domain=-3:5,purple,thick](h)
  \tkzFnDraw[domain=-3:5,red,dashed](f)
  \tkzFnDraw[domain=-3:0,samples=2,{[-)},thick]{2}
  \tkzFnDraw[domain=0:5,samples=2,{[-)},thick]{x/2}
\end{tikzpicture}

\end{document}
